\vspace{-5mm}
\section{Related Works}
\vspace{-4mm}
\label{sec:related_work}



% Atomistic systems can be viewed as graphs, and graph neural networks (GNNs) have the potential to accelerate quantum mechanics calculations by learning their efficient approximations. }
% There are two main ways to represent atomistic graphs~\cite{atom3d}, which are chemical bond graphs and 3D graphs.
% Chemical bond graphs use edges to represent covalent bonds but do not consider the embedded 3D geometry.



%Here, we focus on equivariant neural networks and leave discussion on other works for the appendix.



%Here, we focus on equivariant neural networks and discuss other works in Sec.~\ref{appendix:sec:related_works} in appendix.
\revision{We focus on equivariant neural networks here. We provide a detailed comparison between other equivariant Transformers and Equiformer and discuss other related works in Sec.~\ref{appendix:sec:related_works} in appendix.}
\vspace{-1.5mm}
\paragraph{\textit{SE(3)/E(3)}-Equivariant GNNs.}
\vspace{-2.5mm}
%A different kind of approaches to modeling 3D graphs is to incorporate equivariant geometric tensors such as type-$L$ vectors as discussed in Sec.~\ref{subsec:irreducible_representations} into features of GNNs~\cite{tfn, kondor2018clebsch, 3dsteerable, se3_transformer, geometric_prediction, nequip, gvp, painn, egnn, segnn, torchmd_net, eqgat, allergo}.
Equivariant neural networks~\citep{tfn, kondor2018clebsch, 3dsteerable, se3_transformer, l1net, geometric_prediction, nequip, gvp, painn, egnn, se3_wavefunction, segnn, torchmd_net, eqgat, allergo} operate on geometric tensors like type-$L$ vectors to achieve equivariance.
The central idea is to use functions of geometry built from spherical harmonics and irreps features to achieve 3D rotational and translational equivariance as proposed in Tensor Field Network (TFN)~\citep{tfn}, which generalizes 2D counterparts~\citep{harmonis_network, group_equivariant_convolution_networks, spherical_cnn} to 3D Euclidean space~\citep{tfn, 3dsteerable, kondor2018clebsch}.
Previous works differ in equivariant operations used in their networks.
TFN~\citep{tfn} and NequIP~\citep{nequip} use graph convolution with linear messages, with the latter utilizing extra equivariant gate activations~\citep{3dsteerable}.
SEGNN~\citep{segnn} introduces non-linear messages~\citep{neural_message_passing_quantum_chemistry, gns} for irreps features, and the non-linear messages use the same gate activation and improve upon linear messages. 
SE(3)-Transformer~\citep{se3_transformer} adopts an equivariant version of dot product (DP)  attention~\citep{transformer} with linear messages, and the attention can support vectors of any type $L$.
Subsequent works on equivariant Transformers~\citep{torchmd_net, eqgat} follow the practice of DP attention and linear messages but use more specialized architectures considering only type-$0$ and type-$1$ vectors.

%SEGNN~\cite{segnn} utilizes non-linear message passing networks~\cite{neural_message_passing_quantum_chemistry, gns} for irreps features.
\vspace{-1.5mm}
The proposed Equiformer incorporates all the advantages through combining MLP attention with non-linear messages and supporting vectors of any type.
Compared to TFN, NequIP, SEGNN and SE(3)-Transformer, the proposed combination of MLP attention and non-linear messages is more expressive than pure linear or non-linear messages and pure MLP or dot product attention.
Compared to other equivariant Transformers~\citep{torchmd_net, eqgat}, in addition to being more expressive, the proposed attention mechanism can support vectors of higher degrees (types) and involve higher order tensor product interactions, which can lead to better performance. %~\citep{nequip, segnn}.

%Besides, Equiformer adopts the same network design as typical Transformer~\cite{layer_norm_in_transformer, vit} by following how modules are connected.

%\todo{Add reference for message passing networks.}


\begin{comment}
\vspace{-2mm}
\subsection{Invariant GNNs}
\vspace{-2mm}
Previous works~\cite{schnet, cgcnn, physnet, dimenet, dimenet_pp, orbnet, spherenet, spinconv, gemnet} extract invariant information from 3D atomistic graphs and operate on the resulting invariant graphs.
They mainly differ in leveraging different geometric information such as distances, bond angles (3 atom features) or dihedral angles (4 atom features).
SchNet~\cite{schnet} uses relative distances and proposes continuous-filter convolutional layers to learn local interaction between atom pairs. 
DimeNet series~\cite{dimenet, dimenet_pp} incorporate bond angles by using triplet representations of atoms.
%SphereNet~\cite{spherenet} and GemNet~\cite{gemnet} further extend to consider dihedral angles, which has been shown by GemNet~\cite{gemnet} to be universal approximators. 
SphereNet~\cite{spherenet} and GemNet~\cite{gemnet, gemnet_oc} further extend to consider dihedral angles for better performance.
In order to consider directional information contained in angles, they rely on triplet or quadruplet representations of atoms.
In addition to being memory-intensive~\cite{gemnet_xl}, they also change graph structures by introducing higher-order interaction terms~\cite{line_graph}, which would require non-trivial modifications to generic GNNs in order to apply them to 3D graphs.
%In contrast, the proposed Equiformer uses equivariant irreps features to consider directional information without complicating graph structures and therefore can directly inherit the design of generic GNNs.
In contrast, the proposed Equiformer uses irreps features to consider directional information without complicating graph structures and therefore can directly inherit the design of generic GNNs.


\vspace{-2mm}
\subsection{Attention and Transformer}
\vspace{-2mm}

\vspace{-1mm}
\paragraph{Graph Attention.}
Graph attention networks (GAT)~\cite{gat, gatv2} use multi-layer perceptrons (MLP) to calculate attention weights in a similar manner to message passing networks.
Subsequent works using graph attention mechanisms follow either GAT-like MLP attention~\cite{relational_graph_attention_networks, graph_attention_with_self_supervision} or Transformer-like dot product attention~\cite{gaan, hard_graph_attention, masked_label_prediction, generalization_transformer_graphs, graph_attention_with_self_supervision, spectral_attention}.
In particular, Kim \textit{et al.}~\cite{graph_attention_with_self_supervision} compares these two types of attention mechanisms empirically under a self-supervised setting.
Brody \textit{et al.}~\cite{gatv2} analyzes their theoretical differences and compares their performance in general settings.
%Since GAT-like MLP attention includes non-linearity when calculating attention weights, from theoretical perspectives, it is more expressive than Transformer-like dot product attention, where only linear operations are used for attention weights.

\vspace{-2mm}
\vspace{-1mm}
\paragraph{Graph Transformer.}
A different line of research focuses on adapting standard Transformer networks to graph problems~\cite{generalization_transformer_graphs, grove, spectral_attention, graphormer, graphormer_3d}.
They adopt dot product attention in Transformers~\cite{transformer} and propose different approaches to incorporate graph-related inductive biases into their networks.
GROVE~\cite{grove} includes additional message passing layers or graph convolutional layers to incorporate local graph structures when calculating attention weights.
SAN~\cite{spectral_attention} proposes to learn position embeddings of nodes with full Laplacian spectrum.
Graphormer~\cite{graphormer} proposes to encode degree information in centrality embeddings and encode distances and edge features in attention biases.
The proposed Equiformer belongs to one of these attempts to generalize standard Transformers to graphs and is dedicated to 3D graphs.
To incorporate 3D-related inductive biases, we adopt an equivariant version of Transformers with irreps features and propose novel equivariant graph attention.

\end{comment}